\QHeading{第1题　口算。}
\par\noindent


\begin{oralgrid}
\OralItem{1}{0.45+0.55}{1}
\OralItem{2}{6.3+0.07}{6.37}
\OralItem{3}{8.9-0.09}{8.81}
\OralItem{4}{10.25-3.5}{6.75}
\OralItem{5}{0.68+0.32}{1}
\OralItem{6}{7.25-0.25}{7}
\OralItem{7}{0.72-0.2}{0.52}
\OralItem{8}{2.45-1.4}{1.05}
\OralItem{9}{0.9+1.2}{2.1}
\OralItem{10}{0.7+2.5}{3.2}
\OralItem{11}{1.8+3.6}{5.4}
\OralItem{12}{5.6-2.9}{2.7}\end{oralgrid}
\par

\QHeading{口算思路补充（对应上表题号）}

\Answer{上面列出得数，下面按小题给出不同计算思路。计数单位法适用于小数，整数题可比较凑整、拆分或商不变等方法。}
\MethodTitle{第1小题：$0.45+0.55$}

\MethodBlock{口算思路一：按相同计数单位计算}{先看成若干个百分之一，按整数加减，最后还原成小数。}{\begin{aligned}[t]&\phantom{={}}0.45+0.55\\
&=(45+55)\div 100\\
&=100\div 100\\
&=1\end{aligned}}
\MethodTitle{第2小题：$6.3+0.07$}

\MethodBlock{口算思路一：按相同计数单位计算}{先看成若干个百分之一，按整数加减，最后还原成小数。}{\begin{aligned}[t]&\phantom{={}}6.3+0.07\\
&=(630+7)\div 100\\
&=637\div 100\\
&=6.37\end{aligned}}
\MethodTitle{第3小题：$8.9-0.09$}

\MethodBlock{口算思路一：按相同计数单位计算}{先看成若干个百分之一，按整数加减，最后还原成小数。}{\begin{aligned}[t]&\phantom{={}}8.9-0.09\\
&=(890-9)\div 100\\
&=881\div 100\\
&=8.81\end{aligned}}

\MethodBlock{口算思路二：补差法}{从减数先补到一个整数，再补到被减数，两段差相加。}{\begin{aligned}[t]&\phantom{={}}8.9-0.09\\
&=(1-0.09)+(8.9-1)\\
&=0.91+7.9\\
&=8.81\end{aligned}}
\MethodTitle{第4小题：$10.25-3.5$}

\MethodBlock{口算思路一：按相同计数单位计算}{先看成若干个百分之一，按整数加减，最后还原成小数。}{\begin{aligned}[t]&\phantom{={}}10.25-3.5\\
&=(1025-350)\div 100\\
&=675\div 100\\
&=6.75\end{aligned}}

\MethodBlock{口算思路二：补差法}{从减数先补到一个整数，再补到被减数，两段差相加。}{\begin{aligned}[t]&\phantom{={}}10.25-3.5\\
&=(4-3.5)+(10.25-4)\\
&=0.5+6.25\\
&=6.75\end{aligned}}
\MethodTitle{第5小题：$0.68+0.32$}

\MethodBlock{口算思路一：按相同计数单位计算}{先看成若干个百分之一，按整数加减，最后还原成小数。}{\begin{aligned}[t]&\phantom{={}}0.68+0.32\\
&=(68+32)\div 100\\
&=100\div 100\\
&=1\end{aligned}}
\MethodTitle{第6小题：$7.25-0.25$}

\MethodBlock{口算思路一：按相同计数单位计算}{先看成若干个百分之一，按整数加减，最后还原成小数。}{\begin{aligned}[t]&\phantom{={}}7.25-0.25\\
&=(725-25)\div 100\\
&=700\div 100\\
&=7\end{aligned}}

\MethodBlock{口算思路二：补差法}{从减数先补到一个整数，再补到被减数，两段差相加。}{\begin{aligned}[t]&\phantom{={}}7.25-0.25\\
&=(1-0.25)+(7.25-1)\\
&=0.75+6.25\\
&=7\end{aligned}}
\MethodTitle{第7小题：$0.72-0.2$}

\MethodBlock{口算思路一：按相同计数单位计算}{先看成若干个百分之一，按整数加减，最后还原成小数。}{\begin{aligned}[t]&\phantom{={}}0.72-0.2\\
&=(72-20)\div 100\\
&=52\div 100\\
&=0.52\end{aligned}}
\MethodTitle{第8小题：$2.45-1.4$}

\MethodBlock{口算思路一：按相同计数单位计算}{先看成若干个百分之一，按整数加减，最后还原成小数。}{\begin{aligned}[t]&\phantom{={}}2.45-1.4\\
&=(245-140)\div 100\\
&=105\div 100\\
&=1.05\end{aligned}}

\MethodBlock{口算思路二：补差法}{从减数先补到一个整数，再补到被减数，两段差相加。}{\begin{aligned}[t]&\phantom{={}}2.45-1.4\\
&=(2-1.4)+(2.45-2)\\
&=0.6+0.45\\
&=1.05\end{aligned}}
\MethodTitle{第9小题：$0.9+1.2$}

\MethodBlock{口算思路一：按相同计数单位计算}{先看成若干个十分之一，按整数加减，最后还原成小数。}{\begin{aligned}[t]&\phantom{={}}0.9+1.2\\
&=(9+12)\div 10\\
&=21\div 10\\
&=2.1\end{aligned}}
\MethodTitle{第10小题：$0.7+2.5$}

\MethodBlock{口算思路一：按相同计数单位计算}{先看成若干个十分之一，按整数加减，最后还原成小数。}{\begin{aligned}[t]&\phantom{={}}0.7+2.5\\
&=(7+25)\div 10\\
&=32\div 10\\
&=3.2\end{aligned}}
\MethodTitle{第11小题：$1.8+3.6$}

\MethodBlock{口算思路一：按相同计数单位计算}{先看成若干个十分之一，按整数加减，最后还原成小数。}{\begin{aligned}[t]&\phantom{={}}1.8+3.6\\
&=(18+36)\div 10\\
&=54\div 10\\
&=5.4\end{aligned}}
\MethodTitle{第12小题：$5.6-2.9$}

\MethodBlock{口算思路一：按相同计数单位计算}{先看成若干个十分之一，按整数加减，最后还原成小数。}{\begin{aligned}[t]&\phantom{={}}5.6-2.9\\
&=(56-29)\div 10\\
&=27\div 10\\
&=2.7\end{aligned}}

\MethodBlock{口算思路二：补差法}{从减数先补到一个整数，再补到被减数，两段差相加。}{\begin{aligned}[t]&\phantom{={}}5.6-2.9\\
&=(3-2.9)+(5.6-3)\\
&=0.1+2.6\\
&=2.7\end{aligned}}


\par\noindent

\QHeading{第2题　脱式计算，能简算的要简算。}

\CalcBlock{1}{\begin{aligned}[t]&\phantom{={}}\color{black}2.56+3.44+5.89\\&\color{answerred}=(2.56+3.44)+5.89\\&\color{answerred}=6+5.89\\&\color{answerred}=11.89\end{aligned}}
\par

\CalcBlock[方法一]{2}{\begin{aligned}[t]&\phantom{={}}\color{black}18.7-5.23-4.77\\&\color{answerred}=18.7-(5.23+4.77)\\&\color{answerred}=18.7-10\\&\color{answerred}=8.7\end{aligned}}
\MethodBlock{方法二：先处理括号或乘除，再多减再加回}{减去凑成的整数后，把多减的部分加回来。}{\begin{aligned}[t]&\phantom{={}}18.7-5.23-4.77\\
&=13.47-4.77\\
&=13.47-5+0.23\\
&=8.47+0.23\\
&=8.7\end{aligned}}

\CalcBlock{3}{\begin{aligned}[t]&\phantom{={}}\color{black}35.6+12.4-20.8\\&\color{answerred}=48-20.8\\&\color{answerred}=27.2\end{aligned}}
\CalcBlock[方法一]{4}{\begin{aligned}[t]&\phantom{={}}\color{black}42.9-16.5+7.1\\&\color{answerred}=42.9+7.1-16.5\\&\color{answerred}=(42.9+7.1)-16.5\\&\color{answerred}=50-16.5\\&\color{answerred}=33.5\end{aligned}}
\MethodBlock{方法二：按运算顺序直接计算（对照）}{本法只作结果核对；本题要求简便计算，前面的简便解法仍保留。}{\begin{aligned}[t]&\phantom{={}}42.9-16.5+7.1\\
&=26.4+7.1\\
&=33.5\end{aligned}}

\CalcBlock{5}{\begin{aligned}[t]&\phantom{={}}\color{black}0.85+1.23+2.92\\&\color{answerred}=2.08+2.92\\&\color{answerred}=5\end{aligned}}
\CalcBlock[方法一]{6}{\begin{aligned}[t]&\phantom{={}}\color{black}60.5-28.36-11.64\\&\color{answerred}=60.5-(28.36+11.64)\\&\color{answerred}=60.5-40\\&\color{answerred}=20.5\end{aligned}}
\MethodBlock{方法二：按运算顺序直接计算（对照）}{本法只作结果核对；本题要求简便计算，前面的简便解法仍保留。}{\begin{aligned}[t]&\phantom{={}}60.5-28.36-11.64\\
&=32.14-11.64\\
&=20.5\end{aligned}}

\par\noindent

\QHeading{第3题　解决问题。}

\Problem{① 一根完整的彩带长25.67米，第一次剪掉7.15米用来绑礼盒，第二次剪掉4.85米用来装饰贺卡。这根彩带被剪短了多少米？还剩多少米？}
\par


\MethodTitle{方法一}


\MethodTitle{方法二：先求剩余，再用总长检验剪短量}
\Answer{按剪掉的先后顺序求剩余；总长减剩余也能得到剪短的长度。}
\Answer{\textbf{剩余长度}\par \Formula{\begin{aligned}[t]&\phantom{={}}25.67-7.15-4.85\\
&=18.52-4.85\\
&=13.67\text{（米）}\end{aligned}}}
\Answer{\textbf{剪短长度}\par \Formula{\begin{aligned}[t]&\phantom{={}}25.67-13.67\\
&=12\text{（米）}\end{aligned}}}
\Answer{答：彩带被剪短了12米。\par 答：还剩13.67米。}
\Problem{② 体育课上，小明立定跳远跳了1.5米，小红比小明少跳0.25米，小亮比小红多跳0.18米。小亮跳了多少米？三人立定跳远的总长度一共是多少米？}

\MethodTitle{方法一}


\MethodTitle{方法二：先比较小亮与小明，再用相同基数求总长}
\Answer{小亮比小明少0.25米又多0.18米，合起来少0.07米；总长度先按三个1.5米计算，再扣除两人的不足。}
\Answer{\textbf{小亮比小明少跳}\par \Formula{\begin{aligned}[t]&\phantom{={}}0.25-0.18\\
&=0.07\text{（米）}\end{aligned}}}
\Answer{\textbf{小亮跳远长度}\par \Formula{\begin{aligned}[t]&\phantom{={}}1.5-0.07\\
&=1.43\text{（米）}\end{aligned}}}
\Answer{\textbf{三人总长度}\par \Formula{\begin{aligned}[t]&\phantom{={}}1.5\times 3-0.25-0.07\\
&=4.5-0.25-0.07\\
&=4.25-0.07\\
&=4.18\text{（米）}\end{aligned}}}
\Answer{答：小亮跳了1.43米。\par 答：三人跳远总长度是4.18米。}


\Answer{
\QHeading{计算过程：3. 解决问题}
\begin{enumerate}
\item 两次剪去的长度为$7.15+4.85=12$（米）。

答：这根彩带被剪短了12米。

剩余长度用综合算式计算：
\Formula{\begin{aligned}[t]&\phantom{={}}25.67-(7.15+4.85)\\&=25.67-12\\&=13.67\text{（米）}\end{aligned}}
答：这根彩带还剩13.67米。
\item 小亮的成绩比小红多0.18米，小红比1.5米少0.25米：
\Formula{\begin{aligned}[t]&\phantom{={}}1.5-0.25+0.18\\&=1.25+0.18\\&=1.43\text{（米）}\end{aligned}}
答：小亮跳了1.43米。

三人的总长度：
\Formula{\begin{aligned}[t]&\phantom{={}}1.5+1.25+1.43\\&=2.75+1.43\\&=4.18\text{（米）}\end{aligned}}
答：三人立定跳远的总长度一共是4.18米。
\end{enumerate}


}
